\documentclass[11pt,a4paper]{article}

\usepackage[utf8]{inputenc}
\usepackage[T1]{fontenc}
\usepackage[french]{babel}
\usepackage[a4paper,margin=2.2cm]{geometry}
\usepackage{array}
\usepackage{booktabs}
\usepackage{longtable}
\usepackage{tabularx}
\usepackage{hyperref}
\usepackage{xcolor}
\usepackage{enumitem}

\hypersetup{
  colorlinks=true,
  linkcolor=black,
  urlcolor=blue
}

\title{Sprint 1 Review\\Marketplace Locale Intelligente}
\author{Thomas -- Lead Developer et Product Owner}
\date{Mars 2026}

\begin{document}

\maketitle

\section{Objet de la review}

Cette revue de sprint a pour objectif de présenter au Product Owner et aux parties prenantes les livrables effectivement intégrés à la fin du sprint 1, d'évaluer leur conformité par rapport au backlog initial et d'identifier les écarts à traiter au sprint 2.

Le sprint 1 était consacré aux fondations du produit: socle backend, base de données PostgreSQL, premiers contrats d'API, initialisation du frontend, premières pratiques DevOps et mise à disposition d'un jeu de données mock pour l'équipe data.

\section{Rappel du périmètre prévu}

\subsection{Objectif sprint}

\begin{quote}
Mettre en place l'architecture de base de la marketplace, sécuriser les premiers flux utilisateur et disposer d'un socle technique suffisamment propre pour enchaîner sur la logique métier au sprint 2.
\end{quote}

\subsection{User stories ciblées}

\begin{itemize}[leftmargin=1.5em]
  \item \textbf{US1.1} En tant que visiteur, je peux créer un compte client ou producteur.
  \item \textbf{US1.2} En tant qu'utilisateur connecté, je peux me connecter et récupérer mon profil.
  \item \textbf{US1.3} En tant que client, je peux consulter la liste des produits et le détail d'un produit.
\end{itemize}

\subsection{Features techniques ciblées}

\begin{itemize}[leftmargin=1.5em]
  \item \textbf{FEAT1.1} Initialisation backend + BDD PostgreSQL.
  \item \textbf{FEAT1.2} CRUD minimal API produits + endpoints auth/utilisateur.
  \item \textbf{FEAT1.3} Initialisation frontend et architecture composants.
  \item \textbf{FEAT1.4} Setup GitFlow + protection de branches + premiers checks CI.
  \item \textbf{FEAT1.5} Premiers Dockerfile et environnement local conteneurisé.
  \item \textbf{FEAT1.6} Génération d'un mock dataset exploitable par l'équipe data.
\end{itemize}

\section{Livrables intégrés au dépôt}

\subsection{Backend et base de données}

Les fondations backend sont en place autour de \textbf{FastAPI} avec un routage versionné \texttt{/api/v1}. Les éléments notables livrés sont:

\begin{itemize}[leftmargin=1.5em]
  \item initialisation de l'application et structuration modulaire des routeurs;
  \item configuration de la couche base de données et modèles SQLAlchemy;
  \item schéma PostgreSQL versionné couvrant utilisateurs, producteurs, produits, commandes, paiements et vue analytique;
  \item dépendances d'authentification et contrats de schémas API.
\end{itemize}

Les endpoints effectivement présents pour le périmètre sprint 1 sont:

\begin{itemize}[leftmargin=1.5em]
  \item \texttt{POST /api/v1/auth/register}
  \item \texttt{POST /api/v1/auth/login}
  \item \texttt{GET /api/v1/users/me}
  \item \texttt{GET /api/v1/products/}
  \item \texttt{GET /api/v1/products/\{id\}}
  \item \texttt{GET /api/v1/producers}
  \item \texttt{GET /api/v1/producers/\{id\}/products}
\end{itemize}

Au-delà du strict sprint 1, le dépôt anticipe déjà des routeurs \texttt{orders} et \texttt{data}, ce qui constitue une avance de structuration utile pour le sprint 2 et pour l'intégration future de la dimension analytique.

\subsection{Frontend}

Le frontend a été initialisé sur \textbf{React + Vite}, avec une base UI cohérente et immédiatement exploitable. Le dépôt contient:

\begin{itemize}[leftmargin=1.5em]
  \item une architecture claire par pages, composants, contexte d'authentification, contexte panier et clients API;
  \item une page d'accueil produit;
  \item les pages \emph{login} et \emph{register};
  \item un catalogue de produits avec recherche locale et filtrage par catégorie;
  \item une page de listing des producteurs;
  \item les premières pages panier et commandes, qui dépassent le strict périmètre du sprint 1.
\end{itemize}

Cette base permet d'enchaîner rapidement sur le tunnel d'achat du sprint 2 sans refonte structurelle.

\subsection{DevOps, CI et conteneurisation}

Les premiers livrables d'industrialisation sont présents:

\begin{itemize}[leftmargin=1.5em]
  \item \textbf{Dockerfile backend} pour exécuter l'API FastAPI dans un conteneur Python 3.11;
  \item \textbf{Dockerfile frontend} multi-stage avec build Vite puis exposition via Nginx;
  \item \textbf{docker-compose} orchestrant PostgreSQL, backend et frontend;
  \item workflow \textbf{CI} déclenché sur \texttt{develop} et \texttt{main};
  \item documentation GitFlow et template de Pull Request.
\end{itemize}

Le pipeline actuellement versionné couvre surtout le backend: compilation Python, test d'import de l'application, exécution de tests backend et build de l'image Docker backend.

\subsection{Volet data}

Le sprint 1 a correctement préparé le terrain pour l'équipe data:

\begin{itemize}[leftmargin=1.5em]
  \item schéma de base pensé à la fois pour le transactionnel et l'analytique;
  \item script de \emph{seed} pour générer un volume significatif de données mock;
  \item notebook d'exploration initial;
  \item documentation dédiée au dataset.
\end{itemize}

Le script de génération prévoit notamment des utilisateurs, producteurs, produits, commandes et paiements simulés, avec des volumes suffisants pour amorcer l'exploration BI et ML sans attendre un trafic réel.

\section{Synthèse par user story et feature}

\begin{longtable}{>{\raggedright\arraybackslash}p{2.2cm} >{\raggedright\arraybackslash}p{2.5cm} >{\raggedright\arraybackslash}p{8.4cm} >{\raggedright\arraybackslash}p{2.3cm}}
\toprule
\textbf{ID} & \textbf{Statut} & \textbf{Justification} & \textbf{Décision PO} \\
\midrule
\endhead
US1.1 & Livrée & Inscription client/producteur disponible via l'API et formulaire frontend dédié. & Acceptée \\
US1.2 & Livrée & Login JWT et récupération du profil implémentés côté backend; gestion du token et du profil présente côté frontend. & Acceptée \\
US1.3 & Partiellement livrée & Listing catalogue disponible côté backend et frontend; endpoint détail produit présent côté API, mais la consultation détaillée n'est pas encore matérialisée par une page frontend dédiée. & Acceptation partielle \\
FEAT1.1 & Livrée & Backend FastAPI initialisé, modèle de données posé, schéma PostgreSQL versionné. & Validée \\
FEAT1.2 & Livrée & Endpoints auth, profil, produits et producteurs disponibles. & Validée \\
FEAT1.3 & Livrée & Frontend React/Vite initialisé avec UI kit et structure applicative propre. & Validée \\
FEAT1.4 & Partiellement livrée & GitFlow documenté, CI backend présente, mais la couverture qualité reste limitée et la protection de branches n'est pas démontrable depuis le dépôt seul. & Validation partielle \\
FEAT1.5 & Livrée & Dockerfiles front/back et environnement local conteneurisé disponibles. & Validée \\
FEAT1.6 & Livrée & Mock dataset, notebook d'exploration et base orientée analytics disponibles. & Validée \\
\bottomrule
\end{longtable}

\section{Démonstration fonctionnelle présentable}

À l'issue du sprint 1, la démonstration de review peut s'articuler autour du scénario suivant:

\begin{enumerate}[leftmargin=1.5em]
  \item présentation de l'architecture cible et de la séparation frontend / backend / base de données;
  \item démonstration de l'inscription d'un utilisateur client ou producteur;
  \item démonstration de l'authentification et de la récupération du profil utilisateur;
  \item consultation du catalogue et filtrage local des produits;
  \item mise en avant de la page producteurs;
  \item ouverture du schéma SQL et du script de seed pour illustrer la préparation du volet data;
  \item présentation du \texttt{docker-compose.yml} et du workflow CI pour montrer la base DevOps du projet.
\end{enumerate}

\section{Points forts du sprint}

\begin{itemize}[leftmargin=1.5em]
  \item \textbf{Sprint utile et structurant}: la plupart des fondations critiques ont été intégrées tôt.
  \item \textbf{Bonne cohérence d'ensemble}: le dépôt relie déjà backend, base, frontend, Docker et data.
  \item \textbf{Anticipation produit}: certaines briques du sprint 2 sont déjà préparées dans l'architecture.
  \item \textbf{Vision data présente dès le départ}: le schéma et les données mock ne sont pas ajoutés tardivement, ce qui réduit le risque de dette analytique.
\end{itemize}

\section{Écarts, risques et limites observés}

\begin{itemize}[leftmargin=1.5em]
  \item le détail produit n'est pas encore entièrement visible côté frontend, ce qui empêche de considérer \textbf{US1.3} comme totalement terminée;
  \item la couverture qualité est encore modeste: peu de tests versionnés et pipeline CI centré sur le backend uniquement;
  \item certaines routes frontend semblent préparées plus vite que finalisées, ce qui peut générer des redirections vers des écrans non encore implémentés;
  \item la validation locale complète n'a pas pu être rejouée dans l'environnement courant sans installation préalable des dépendances Python et Node.
\end{itemize}

\section{Décisions de review}

En tant que Product Owner, je formule les décisions suivantes:

\begin{itemize}[leftmargin=1.5em]
  \item \textbf{Acceptation} de l'incrément de fondation backend, base, data et conteneurisation;
  \item \textbf{Acceptation} du socle frontend comme base de travail pour la suite du produit;
  \item \textbf{Acceptation partielle} de la couverture fonctionnelle utilisateur, avec report explicite du détail produit complet au début du sprint 2 si nécessaire;
  \item \textbf{Priorité élevée} au sprint 2 sur la consolidation de l'intégration front-back, l'alignement des routes frontend et l'élargissement de la validation automatisée.
\end{itemize}

\section{Plan de passage au sprint 2}

Les priorités de transition sont les suivantes:

\begin{enumerate}[leftmargin=1.5em]
  \item finaliser les écarts de finition du sprint 1, en particulier sur les écrans ou routes incomplètes;
  \item développer le cycle panier $\rightarrow$ commande $\rightarrow$ paiement simulé;
  \item renforcer l'intégration continue avec un build frontend et davantage de tests automatisés;
  \item conserver la discipline GitFlow et la revue de code croisée sur les branches \texttt{feature/*}.
\end{enumerate}

\section{Conclusion}

Le sprint 1 est globalement \textbf{réussi}. Les fondations essentielles sont présentes, la base technique est crédible et le projet est dans une position saine pour attaquer le coeur métier au sprint 2. Le principal enjeu n'est plus de poser l'architecture, mais de consolider la qualité d'intégration et d'achever quelques parcours front pour transformer ce socle en expérience utilisateur pleinement cohérente.

\end{document}